\documentclass[11pt]{article}
\usepackage[margin=1in]{geometry}
\usepackage[T1]{fontenc}
\usepackage[utf8]{inputenc}
\usepackage{lmodern}
\usepackage{graphicx}
\usepackage{booktabs}
\usepackage{tabularx}
\usepackage{longtable}
\usepackage{calc}
\usepackage{float}
\usepackage{xcolor}
\usepackage{amsmath,amssymb}
\usepackage[numbers,sort&compress]{natbib}
\usepackage[hidelinks]{hyperref}
\providecommand{\doi}[1]{\href{https://doi.org/#1}{https://doi.org/#1}}
\setlength{\emergencystretch}{3em}
\setlength{\parindent}{1.5em}
\setlength{\parskip}{0.25em}
\linespread{1.18}
\sloppy
\widowpenalty=10000
\clubpenalty=10000
\renewcommand{\refname}{References}
\providecommand{\tightlist}{\setlength{\itemsep}{0pt}\setlength{\parskip}{0pt}}
\providecommand{\real}[1]{#1}
\title{Supplementary Information: Operator identity and out-of-user generalization in automated surgical skill assessment: a multi-dataset evaluation}
\author{Hamilton Jeong\textsuperscript{1,*}, Zaid Elnasser\textsuperscript{1},
Laurienne Jeong\textsuperscript{1},\\
Su\'elia de Siqueira Rodrigues Fleury Rosa\textsuperscript{2}\\[0.5em]
\small International Journal of Computer Assisted Radiology and Surgery\\
\small \textsuperscript{1}College of Agriculture and Life Sciences, Cornell University, Ithaca, NY, USA\\
\small \textsuperscript{2}Meinig School of Biomedical Engineering, Cornell University, Ithaca, NY, USA\\
\small \textsuperscript{*}Corresponding author: \href{mailto:hhj22@cornell.edu}{hhj22@cornell.edu}}
\date{}
\begin{document}
\maketitle

\setcounter{secnumdepth}{0}
\section{S1. Purpose and locked analysis
boundary}\label{s1.-purpose-and-locked-analysis-boundary}

This supplement documents the study population, outcome provenance,
preprocessing, folds, models, interventions, inference, robustness
analyses, deviations, and reproducibility controls. It accompanies the
original Phase 5 locked outputs and the post-lock correction of 16
November 2026 (S15). Primary outcome and intervention models were not
refit. A separate post-lock fixed-C identity diagnostic is documented in
S16. The exact private out-of-fold inputs were verified against
\texttt{data/manifests/phase3\_artifacts.csv} and
\texttt{data/manifests/phase4\_artifacts.csv}; Phase 5 derived all
inferential summaries from those files.

The protocol was frozen as Git tag \texttt{protocol-v1} on 21 September
2026. The Phase 5 artifact manifest is
\texttt{data/manifests/phase5\_artifacts.csv}, whose locked manifest
SHA-256 is
\nolinkurl{a4345fc8cdcb50b90e286914ff6dec8ec1c7373e248f60a260c4eaf11f0d9bc3}.

\section{S2. Dataset acquisition, licenses, and
cohorts}\label{s2.-dataset-acquisition-licenses-and-cohorts}

\subsection{S2.1 JIGSAWS}\label{s2.1-jigsaws}

The JHU-ISI Gesture and Skill Assessment Working Set contains
synchronized kinematics, video, task transcriptions, modified-OSATS
ratings, and self-declared experience groups \citep{Ahmidi2017}. Access
is governed by an academic research agreement. An authorized local copy
was connected by a Git-ignored symbolic link; raw or participant-level
derived records were not committed. We used 36 knot-tying, 28
needle-passing, and 39 suturing trials after exact matching of
kinematics, transcriptions, and ratings. Needle passing represented
seven participants; the other tasks represented eight.

Total GRS was verified as the sum of six rating elements for every
included trial. The novice/intermediate/expert label was treated as
person-constant and was not represented as a trial-varying expert
judgment.

\subsection{S2.2 ROSMA}\label{s2.2-rosma}

ROSMA version \texttt{10.5281/zenodo.7100542} is distributed under CC BY
4.0 and contains kinematics and task scoring for 12 participants
\citep{RivasBlanco2023, ROSMAData2022}. Exact kinematic-score matching
yielded 71 pea-on-a-peg, 65 post-and-sleeve, and 70 wire-chaser-I
trials. The deposited score table had one additional wire-chaser target
without matching kinematics. Observed task errors were retained.
Completion time and the composite source score were excluded because
their formulas use duration, which is directly represented in the
inputs.

Errors were zero inflated: 30/71 pea-on-a-peg, 14/65 post-and-sleeve,
and 17/70 wire-chaser trials had nonzero errors. A deposited
pea-on-a-peg record contained isolated finite coordinate spikes. They
were retained because the protocol prohibited discretionary post-result
exclusion.

\subsection{S2.3 BGU-SKILLS}\label{s2.3-bgu-skills}

BGU-SKILLS version 1.0 (\texttt{10.17617/3.CYSIAP}) is distributed under
CC BY 4.0 \citep{BGUSkillsData2025}. The study enrolled 21 residents and
reported three withdrawals, leaving 18 participants and 972 scheduled
trials. We followed the source inclusion notes, retaining 311
knot-tying, 312 ring-transfer, and 315 suturing trials (938 total).
Completion time, path length, and orientation-change rate are direct
functions of kinematics and were ineligible as independent outcomes. The
dataset therefore contributed to identity recovery and longitudinal
robustness only.

\subsection{S2.4 BGU-PatCut}\label{s2.4-bgu-patcut}

BGU-PatCut version 2.1 (\texttt{10.17617/3.JK4ANF}) is distributed under
CC BY 4.0 \citep{SharonPatCut2024, PatCutData2024}. The cohort comprised
30 force-perturbation and 24 motion-perturbation participants, each
scheduled for 24 trials. Source missing-trial lists yielded 700 and 554
available trials, respectively. Total cutting error
(\texttt{TE\ =\ Eoutside\ +\ Einside}) was verified for each participant
and retained because it was measured from an independent post-trial
gauze image. Task time and the composite efficiency score were excluded
because they contain duration.

\subsection{S2.5 Cohort summary}\label{s2.5-cohort-summary}

{\def\LTcaptype{none} % do not increment counter
\begin{longtable}[]{@{}
  >{\raggedright\arraybackslash}p{(\linewidth - 8\tabcolsep) * \real{0.1765}}
  >{\raggedright\arraybackslash}p{(\linewidth - 8\tabcolsep) * \real{0.1765}}
  >{\raggedleft\arraybackslash}p{(\linewidth - 8\tabcolsep) * \real{0.2353}}
  >{\raggedleft\arraybackslash}p{(\linewidth - 8\tabcolsep) * \real{0.2353}}
  >{\raggedright\arraybackslash}p{(\linewidth - 8\tabcolsep) * \real{0.1765}}@{}}
\toprule\noalign{}
\begin{minipage}[b]{\linewidth}\raggedright
Dataset
\end{minipage} & \begin{minipage}[b]{\linewidth}\raggedright
Tasks
\end{minipage} & \begin{minipage}[b]{\linewidth}\raggedleft
Trials
\end{minipage} & \begin{minipage}[b]{\linewidth}\raggedleft
Participants
\end{minipage} & \begin{minipage}[b]{\linewidth}\raggedright
Performance-analysis role
\end{minipage} \\
\midrule\noalign{}
\endhead
\bottomrule\noalign{}
\endlastfoot
JIGSAWS & Knot tying; needle passing; suturing & 103 & 8 overall &
Expert-rated GRS; person-level mechanism \\
ROSMA & Pea on a peg; post and sleeve; wire chaser I & 206 & 12 &
Objective error counts \\
BGU-SKILLS & Knot tying; ring tower transfer; suturing & 938 & 18 &
Identity and longitudinal analysis only \\
BGU-PatCut & Pattern cutting & 1,254 & 54 & Image-derived total cutting
error \\
\end{longtable}
}

The sum is 2,501 trials and 92 dataset-participants. Cohorts were
distinct; the count does not assert that identity was linked across
datasets.

\section{S3. Target provenance and construct
language}\label{s3.-target-provenance-and-construct-language}

{\def\LTcaptype{none} % do not increment counter
\begin{longtable}[]{@{}
  >{\raggedright\arraybackslash}p{(\linewidth - 6\tabcolsep) * \real{0.2500}}
  >{\raggedright\arraybackslash}p{(\linewidth - 6\tabcolsep) * \real{0.2500}}
  >{\raggedright\arraybackslash}p{(\linewidth - 6\tabcolsep) * \real{0.2500}}
  >{\raggedright\arraybackslash}p{(\linewidth - 6\tabcolsep) * \real{0.2500}}@{}}
\toprule\noalign{}
\begin{minipage}[b]{\linewidth}\raggedright
Dataset
\end{minipage} & \begin{minipage}[b]{\linewidth}\raggedright
Retained target
\end{minipage} & \begin{minipage}[b]{\linewidth}\raggedright
Measurement source
\end{minipage} & \begin{minipage}[b]{\linewidth}\raggedright
Permitted interpretation
\end{minipage} \\
\midrule\noalign{}
\endhead
\bottomrule\noalign{}
\endlastfoot
JIGSAWS & GRS & Sum of six expert modified-OSATS ratings & Expert-rated
trial performance \\
JIGSAWS & experience class & Self-declared label fixed by person &
Person-level mechanism diagnostic \\
ROSMA & task errors & Observed task-specific error/drop count &
Objective task performance \\
BGU-PatCut & total cutting error & Post-trial image deviation &
Objective pattern-cutting accuracy \\
BGU-SKILLS & none & Source metrics derive from kinematics & No
independent performance target \\
\end{longtable}
}

The manuscript does not call ROSMA, BGU-PatCut, or BGU-SKILLS
dataset-native metrics expert skill ratings. It makes no claim about
clinical outcomes, competence, certification, or patient safety.

\section{S4. Deterministic
preprocessing}\label{s4.-deterministic-preprocessing}

Each loader emitted one record per complete trial with a source hash,
participant and trial keys, task, repetition/session/condition metadata,
outcome provenance, task duration, and two timestamped position streams.
Source identifiers were used to build groups but were never model
predictors. Committed fold files use 16-character SHA-256-derived opaque
keys.

Task windows followed source annotations. For each arm, features were
sampling rate; path length; speed mean, standard deviation, and
10th/50th/90th quantiles; mean and 90th-quantile absolute acceleration;
mean and 90th-quantile absolute jerk; idle fraction below 0.001 m/s;
workspace bounding-box volume; and recording-gap fraction. Gripper mean,
standard deviation, range, and total variation were included where
available. Cross-arm features were aligned-speed correlation and
path-length asymmetry. Duration and arm sample counts were included as
direct duration equivalents.

Non-finite rows and duplicate/backward timestamps were removed where
applicable. A displacement spanning more than three median sampling
intervals was excluded from path and derivative calculations.
Interpolated values farther than 2.5 nominal samples from an observation
were excluded from bimanual correlation. Imputation, scaling,
residualization, identity fitting, and hyperparameter selection were
trained within the applicable outer-training fold.

Two end-to-end extractions produced an identical private feature-matrix
hash:
\nolinkurl{0c351007805f52f41459f622faa3c979f0b410685006763fbfcba1fc431a1e50}.

\section{S5. Fold construction and leakage
protection}\label{s5.-fold-construction-and-leakage-protection}

The same-participant manifest assigned complete trials by source
repetition: five folds for JIGSAWS, six for ROSMA and BGU-SKILLS, and 24
for BGU-PatCut. Every participant-task group occupied at least two
folds. The held-out-participant manifest used one fold per participant
for JIGSAWS, ROSMA, and BGU-SKILLS and 10 balanced participant groups
for BGU-PatCut. A participant was indivisible across tasks and trials in
held-out evaluation.

Nested validation matched the intended outer unit. In a
held-out-participant outer fold, inner splits also held out complete
training participants. The outer test fold could not influence
imputation, scaling, ridge penalty, logistic regularization, identity
direction, projection rank, adversarial lambda, or probe
hyperparameters. A leakage canary confirmed that identifier-derived
columns were rejected from production predictors and yielded perfect
identity prediction only when intentionally exposed.

\section{S6. Baseline models and
diagnostics}\label{s6.-baseline-models-and-diagnostics}

Continuous targets used a training-fold mean/median reference,
duration-only ridge, and full-feature ridge. Phase 3 baseline ridge
alpha was selected from \texttt{\{0.1,\ 1,\ 10,\ 100,\ 1000\}}. Phase 4
linear interventions and their I0 comparator instead used fixed
downstream ridge alpha=10; representation selection still used training
data. JIGSAWS skill class used majority, duration-only multinomial
logistic, and full-feature multinomial logistic models; C was selected
from \texttt{\{0.01,\ 0.1,\ 1,\ 10\}}.

Fresh identity probes used multinomial logistic regression and random
forest. Phase 3 forests used 200 trees; Phase 4 used 30 trees to limit
the intervention grid's compute cost, a dated deviation that retained
the same nonlinear model family. Phase 3 ran 50 record-wise identity
permutations per dataset-task, 50 continuous-outcome negative-control
permutations per eligible dataset-task, and 200 paired participant-wise
and record-wise skill permutations per JIGSAWS task.

The JIGSAWS replication compared six skill-accuracy and six GRS-MAE
endpoints with the earlier study (unpublished manuscript; project
repository: https://github.com/hahnyu/surgical-skill). Eleven of 12 were
within frozen tolerances. Held-out-person needle-passing GRS MAE was
3.063 here versus 4.910 previously, a 1.8467-point difference beyond the
1.0-point tolerance. The present pipeline used 41 harmonized variables,
transcription windows, exact repetition folds, fold-local preprocessing,
and nested penalty selection; the earlier pipeline used 16 features at 5
Hz. This discrepancy is reported rather than treated as an independent
replication success.

\section{S7. Intervention algorithms}\label{s7.-intervention-algorithms}

\subsection{S7.1 I1: direct duration
removal}\label{s7.1-i1-direct-duration-removal}

I1 removed \texttt{duration\_s}, \texttt{psm1\_samples}, and
\texttt{psm2\_samples}. All remaining processing matched I0.

\subsection{S7.2 I2: fold-local duration
residualization}\label{s7.2-i2-fold-local-duration-residualization}

Within each training fold, every non-duration feature was regressed on
duration. The fixed coefficients were applied to validation and test
rows. Scaling and prediction remained training fold local.

\subsection{S7.3 I3: linear identity-subspace
removal}\label{s7.3-i3-linear-identity-subspace-removal}

Within each outer-training fold, a regularized multinomial participant
classifier produced a coefficient row space. Candidate ranks were
\texttt{\{0,1,2,4,8,16\}}, clipped to the training fold's maximum
identity rank. For data with a performance target, inner validation
retained ranks with task MAE within 0.05 training-target SD of rank zero
and selected the one with lowest fresh linear identity accuracy,
breaking ties toward smaller rank. For BGU-SKILLS, selection used
identity accuracy alone. The fixed projection was applied to the outer
test fold. This procedure resembles linear guarding by iterative
nullspace projection \citep{Ravfogel2020}, while using the prespecified
coefficient-span implementation.

A random orthogonal projection removed the same fold-specific dimension.
It controlled for capacity loss without targeting identity.

\subsection{S7.4 I4: adversarial summary
representation}\label{s7.4-i4-adversarial-summary-representation}

I4 used a 32-unit hidden layer, 8-unit latent layer, linear task head,
and multinomial identity head. Gradient reversal multiplied the
encoder's identity-loss gradient by negative lambda while the identity
head minimized cross-entropy \citep{Ganin2016}. Candidate lambda values
were \texttt{\{0,0.01,0.1,0.3,1.0\}}. Two participant-safe inner splits
and 30 selection epochs chose the lowest identity-head training accuracy
among candidates with task MAE within 0.05 training-target SD of lambda
zero; final models trained for 80 epochs at learning rate 0.01 and
weight decay 0.0001. Three fixed seeds were used. Controls were the same
architecture at lambda zero and the selected lambda with shuffled
identity labels.

\section{S8. Verification and locked primary
decision}\label{s8.-verification-and-locked-primary-decision}

Representations were frozen before fresh probes were trained under
same-participant folds. Held-out-participant representations were fitted
separately, so these checks do not verify each held-out transformation.
The corrected I4-minus-capacity linear balanced-accuracy difference is
-0.048 (95\% interval -0.098 to +0.002), and the nonlinear difference is
-0.003 (-0.045 to +0.043). Neither verifies reduction. I3-minus-I0
differences are -0.184 (-0.214 to -0.153) for the linear probe and
-0.138 (-0.175 to -0.100) for the nonlinear probe. The original I4-to-I3
decision remains unchanged; these corrected intervals were computed
after the original results were known. I3 aggregates ten dataset-task
groups and I4 seven outcome-eligible groups; the two interventions were
not compared head to head on a common aggregate cohort.

The dimension-matched random projection was worse than I3 by +0.167 on
the linear probe and +0.149 on the nonlinear probe. Duration removal and
residualization did not reproduce I3's joint reduction.

\section{S9. Inference}\label{s9.-inference}

The independent resampling unit was the participant. Participant-average
metrics gave each participant equal weight irrespective of trial count.
Intervals used 2,000 percentile bootstrap replicates within
dataset-task. For corrected aggregate identity verification, seeds were
first averaged within participant and task. Participants were then
resampled within dataset, with each sampled participant carrying all
available tasks, preserving cross-task dependence. Task groups were
averaged equally. The seed was 20261116 and percentile endpoints used
linear interpolation at 0.025 and 0.975. Any draw with no sampled
participant for a task would be rejected in its entirety; no such draws
occurred. Intervals condition on the saved predictions and do not
include model-refitting uncertainty. Leave-one-participant summaries
omit scored participant effects without refitting models.

The held-out-task standardized effect for trial \emph{i} was

\[\frac{|y_i-\widehat{y}_{i,I3}|-|y_i-\widehat{y}_{i,I0}|}{\operatorname{SD}(y_{\mathrm{training,outer\ fold}})}.\]

Negative values favor I3. Participant means were then averaged.
Native-scale effects used the same paired numerator without
standardization. Protocol gaps used the analogous held-out minus
same-participant error after protocol-specific training-fold
standardization.

The archived identity null used fixed-C=1 logistic regression, whereas
observed logistic and forest probes selected hyperparameters. The
formerly reported identity p values are withdrawn because the pipelines
were not matched. No formal permutation-test claim is made for those
descriptive accuracies. The participant-wise versus record-wise JIGSAWS
comparison reports the distribution of paired accuracy differences under
two distinct shuffling schemes. Its plus-one nonpositive tail fraction
is not a valid hypothesis-test p value for a common exchangeable null
and is no longer labeled as one.

\section{S10. Numerical results with corrected
interpretation}\label{s10.-numerical-results-with-corrected-interpretation}

\subsection{S10.1 Unsuppressed identity
recovery}\label{s10.1-unsuppressed-identity-recovery}

{\def\LTcaptype{none} % do not increment counter
\begin{longtable}[]{@{}lrr@{}}
\toprule\noalign{}
Dataset/task & Linear BA & Nonlinear BA \\
\midrule\noalign{}
\endhead
\bottomrule\noalign{}
\endlastfoot
JIGSAWS / Knot Tying & 0.869 & 0.827 \\
JIGSAWS / Needle Passing & 0.943 & 0.840 \\
JIGSAWS / Suturing & 0.944 & 0.844 \\
ROSMA / Pea on a Peg & 0.381 & 0.397 \\
ROSMA / Post and Sleeve & 0.446 & 0.381 \\
ROSMA / Wire Chaser I & 0.228 & 0.431 \\
BGU-SKILLS / Knot Tying & 0.530 & 0.393 \\
BGU-SKILLS / Ring Tower Transfer & 0.596 & 0.487 \\
BGU-SKILLS / Suturing & 0.762 & 0.618 \\
BGU-PatCut / Pattern Cutting & 0.834 & 0.750 \\
\end{longtable}
}

Observed balanced accuracies are descriptive. No matched permutation p
value is available for either tuned model.

\subsection{S10.2 JIGSAWS skill
permutations}\label{s10.2-jigsaws-skill-permutations}

{\def\LTcaptype{none} % do not increment counter
\begin{longtable}[]{@{}
  >{\raggedright\arraybackslash}p{(\linewidth - 8\tabcolsep) * \real{0.1579}}
  >{\raggedleft\arraybackslash}p{(\linewidth - 8\tabcolsep) * \real{0.2105}}
  >{\raggedleft\arraybackslash}p{(\linewidth - 8\tabcolsep) * \real{0.2105}}
  >{\raggedleft\arraybackslash}p{(\linewidth - 8\tabcolsep) * \real{0.2105}}
  >{\raggedleft\arraybackslash}p{(\linewidth - 8\tabcolsep) * \real{0.2105}}@{}}
\toprule\noalign{}
\begin{minipage}[b]{\linewidth}\raggedright
Task
\end{minipage} & \begin{minipage}[b]{\linewidth}\raggedleft
Observed
\end{minipage} & \begin{minipage}[b]{\linewidth}\raggedleft
Participant-wise median {[}95\% range{]}
\end{minipage} & \begin{minipage}[b]{\linewidth}\raggedleft
Record-wise median {[}95\% range{]}
\end{minipage} & \begin{minipage}[b]{\linewidth}\raggedleft
Nonpositive paired differences
\end{minipage} \\
\midrule\noalign{}
\endhead
\bottomrule\noalign{}
\endlastfoot
Knot Tying & 0.917 & 0.778 {[}0.610, 0.917{]} & 0.306 {[}0.167, 0.501{]}
& 0 / 200 \\
Needle Passing & 0.964 & 0.875 {[}0.750, 1.000{]} & 0.321 {[}0.143,
0.501{]} & 0 / 200 \\
Suturing & 0.872 & 0.923 {[}0.846, 0.974{]} & 0.359 {[}0.179, 0.538{]} &
0 / 200 \\
\end{longtable}
}

\subsection{S10.3 Primary held-out-participant
effects}\label{s10.3-primary-held-out-participant-effects}

{\def\LTcaptype{none} % do not increment counter
\begin{longtable}[]{@{}
  >{\raggedright\arraybackslash}p{(\linewidth - 4\tabcolsep) * \real{0.2727}}
  >{\raggedleft\arraybackslash}p{(\linewidth - 4\tabcolsep) * \real{0.3636}}
  >{\raggedleft\arraybackslash}p{(\linewidth - 4\tabcolsep) * \real{0.3636}}@{}}
\toprule\noalign{}
\begin{minipage}[b]{\linewidth}\raggedright
Dataset/task
\end{minipage} & \begin{minipage}[b]{\linewidth}\raggedleft
Standardized I3-I0 MAE {[}95\% interval{]}
\end{minipage} & \begin{minipage}[b]{\linewidth}\raggedleft
Native I3-I0 MAE {[}95\% interval{]}
\end{minipage} \\
\midrule\noalign{}
\endhead
\bottomrule\noalign{}
\endlastfoot
BGU-PatCut / Pattern Cutting & +0.029 {[}-0.038, +0.091{]} & +487.183
{[}-648.849, +1,613.362{]} \\
JIGSAWS / Knot Tying & +0.042 {[}-0.039, +0.132{]} & +0.205 {[}-0.195,
+0.604{]} \\
JIGSAWS / Needle Passing & +0.125 {[}-0.018, +0.290{]} & +0.618
{[}-0.045, +1.400{]} \\
JIGSAWS / Suturing & -0.083 {[}-0.352, +0.170{]} & -0.401 {[}-1.757,
+0.938{]} \\
ROSMA / Pea on a Peg & +24.141 {[}-0.019, +72.407{]} & +28.858
{[}-0.020, +86.539{]} \\
ROSMA / Post and Sleeve & +0.009 {[}-0.125, +0.154{]} & +0.005
{[}-0.072, +0.086{]} \\
ROSMA / Wire Chaser I & +0.004 {[}+0.000, +0.011{]} & +0.003 {[}0.000,
+0.008{]} \\
\end{longtable}
}

The compatible ranges are the study's practical precision statement.
With no validated cross-dataset clinical threshold and no equivalence
margin, intervals crossing zero are not evidence of equivalence. The
pea-on-a-peg standardized result is dominated by small fold-specific
training SDs and should be read with its native result and source
anomaly.

\subsection{S10.4 Same-participant JIGSAWS skill
effects}\label{s10.4-same-participant-jigsaws-skill-effects}

I3 minus I0 participant-average accuracy was -0.025 {[}-0.075, 0.000{]}
for knot tying, 0.000 {[}0.000, 0.000{]} for needle passing, and +0.025
{[}-0.100, +0.125{]} for suturing. The direction was not consistent.

\section{S11. Robustness and influence
analyses}\label{s11.-robustness-and-influence-analyses}

Fold-local I3 selected positive rank in 112/144 outer folds. I4 selected
nonzero lambda in 116/126 folds: 0.01 twice, 0.1 fourteen times, 0.3
forty-six times, and 1.0 fifty-four times; ten folds retained zero.
Across I4 seeds, JIGSAWS skill-accuracy ranges were at most 0.179.
Identity-probe seed ranges had median 0.072 and maximum 0.225.

Participant-level task effects, protocol gaps, identity effects, skill
effects, and leave-one-participant summaries are in
\texttt{results/locked/}. The 748,000 bootstrap replicate rows and 5,226
leave-one rows are preserved in machine-readable files. The supplement
tables in \texttt{results/tables/} document fold composition, selected
hyperparameters, missingness, and seed stability.

\section{S12. Deviations from the prospective
plan}\label{s12.-deviations-from-the-prospective-plan}

\begin{enumerate}
\def\labelenumi{\arabic{enumi}.}
\tightlist
\item
  Phase 0 and Phase 1 were completed in one working session. The
  protocol was still frozen before model fitting.
\item
  Source inspection required explicit non-finite-row and time-gap
  handling. These rules were fixed before modeling and tested
  deterministically.
\item
  The initial same-participant fold draft rotated assignments by
  participant. It was corrected to exact repetition mapping before the
  first Phase 3 prediction.
\item
  I4 used a summary-feature MLP instead of the planned sequence encoder
  because Phase 2 had frozen summary representations rather than
  resampled sequences. This limits claims to summary features.
\item
  Phase 4 nonlinear probes used 30 trees rather than Phase 3's 200 to
  control intervention-grid cost. The nonlinear model family and
  participant-safe tuning remained unchanged.
\item
  Phase 5 regenerated inference from hashed Phase 3/4 predictions rather
  than refitting under a new commit. This preserved the recorded
  generating commit and enabled exact reproducibility.
\item
  Optional BGU component-outcome post hoc models and per-epoch I4 loss
  exports were not produced. Neither was required for the locked primary
  endpoint.
\end{enumerate}

The original deviations did not change participants, primary targets, or
the I4-to-I3 fallback rule. The post-lock correction in S15 changes
aggregate identity uncertainty and inferential reporting after results
were known; it is not presented as a prespecified analysis.

\section{S13. Relationship to the earlier JIGSAWS
study}\label{s13.-relationship-to-the-earlier-jigsaws-study}

The earlier manuscript, ``What does a surgical skill model measure?
Quantifying identity confounding and miscalibration on JIGSAWS,''
analyzed the same 103 JIGSAWS trials (unpublished manuscript; project
repository: https://github.com/hahnyu/surgical-skill). It reported
same-person versus held-person performance, identity recovery, label
permutations, calibration, selective prediction, and a commercial
zero-shot model. The present work reproduces its central
protocol/identity diagnosis using a different 41-feature representation
and exact fold-local pipeline, then extends the question to three
external datasets and directly intervenes on identity representations.
It does not reuse or re-report the commercial model, calibration, or
selective-prediction analyses. The JIGSAWS results are therefore
overlapping replication and extension evidence, not an independent
cohort.

\section{S14. Reproduction}\label{s14.-reproduction}

After authorized data acquisition and environment setup, run:

\begin{verbatim}
.venv/bin/python scripts/run_phase5_locked.py
.venv/bin/python scripts/validate_phase6.py
\end{verbatim}

The Phase 5 command validates Phases 1--4, verifies private artifact
hashes, builds Phase 5 twice, compares every output hash, runs the test
suite, and applies the Phase 5 validator. The Phase 6 validator checks
manuscript structure, bibliography keys, locked numerical statements,
overlap disclosure, construct language, and the claim-evidence/reporting
matrices.

\section{S15. Post-lock technical correction, 16 November
2026}\label{s15.-post-lock-technical-correction-16-november-2026}

The original Phase 3--5 files and their manifest hashes are preserved as
historical artifacts. The current submission supersedes their aggregate
identity intervals and unmatched permutation-test claims. Corrected
sources are in \texttt{results/corrections/2026-11-16/}; the analysis
manifest records input and output hashes, and the package manifest
records submission sources and PDFs.

The correction (1) keeps each dataset-participant's tasks together
during aggregate identity resampling after seed averaging; (2) withdraws
unmatched identity-test p values and relabels the two-scheme JIGSAWS
shuffle comparison descriptively; (3) replaces the pooled cross-task
utility plot with identity effects and explicitly matched controls; (4)
uses participant-average skill accuracy and the Table 5 GRS intervals in
Figure 6; and (5) clarifies fixed versus tuned ridge penalties, the
same-participant scope of fresh probes, and the lack of refitting in
uncertainty and influence calculations. I3 remains the primary
intervention, and no consistent held-out utility benefit emerges.

The corrected bootstrap and figure data use saved aggregate participant
effects and metrics. Private feature and prediction files were initially
cloud-evicted, but became locally available during the correction. The
primary outcome and intervention predictions were unchanged. A separate
matched fixed-C identity diagnostic was then verified as described in
S16; it does not validate the withdrawn p values for the tuned probes.

Reproduce the correction and submission after installing the documented
dependencies:

\begin{verbatim}
python scripts/build_technical_corrections.py
python submission/ijcars/build_submission.py
python scripts/validate_ijcars_submission.py
\end{verbatim}

The historical Phase 5 runner reproduces the archived analysis, not this
correction. The draft validator checks numerical consistency and PDF
rendering diagnostics. The separate \texttt{-\/-final} gate also
requires completed human disclosures, ethics/source-permission
confirmation, author approval, originality confirmation, and public
availability of the corrected analysis release. Outstanding items are
recorded in the submission checklist.

\section{S16. Matched fixed-C identity diagnostic (post-lock exploratory
analysis)}\label{s16.-matched-fixed-c-identity-diagnostic-post-lock-exploratory-analysis}

To address the mismatched original identity null, we fitted an observed
multinomial logistic probe with C=1, matching the original 50
record-wise label-permutation models. Features, trial/fold assignments,
median imputation, standardization, solver, and penalty were identical;
no hyperparameter selection was performed in either pipeline. Feature
and null-prediction hashes were verified against the frozen inputs, and
every permutation preserved the observed trial/fold set and label
counts. This diagnostic was added after inspecting results, rather than
prespecified as the observed test.

Balanced accuracy gives equal weight to identity classes. The unadjusted
Monte Carlo p value is
\texttt{(1\ +\ count(null\ \textgreater{}=\ observed))\ /\ 51}. All ten
observed statistics exceeded all 50 matched null statistics, giving
p=1/51 (approximately 0.01961) in each task. This is the smallest
attainable p value with 50 draws, not evidence for a smaller tail
probability. This is a conditional random-label diagnostic for these
fixed records, not population-level inference about unseen operators.
These exploratory, task-specific p values are not adjusted for multiple
testing and do not establish an omnibus or family-wise significance
claim. They apply only to this fixed-C logistic diagnostic, not to the
tuned logistic or random-forest probes in Table 4.

{\def\LTcaptype{none} % do not increment counter
\begin{longtable}[]{@{}lrr@{}}
\toprule\noalign{}
Dataset/task & Fixed-C balanced accuracy & Matched null median \\
\midrule\noalign{}
\endhead
\bottomrule\noalign{}
\endlastfoot
BGU-PatCut / Pattern Cutting & 0.825 & 0.016 \\
BGU-SKILLS / Knot Tying & 0.530 & 0.051 \\
BGU-SKILLS / Ring Tower Transfer & 0.596 & 0.050 \\
BGU-SKILLS / Suturing & 0.762 & 0.049 \\
JIGSAWS / Knot Tying & 0.869 & 0.084 \\
JIGSAWS / Needle Passing & 0.943 & 0.064 \\
JIGSAWS / Suturing & 0.944 & 0.075 \\
ROSMA / Pea on a Peg & 0.394 & 0.056 \\
ROSMA / Post and Sleeve & 0.504 & 0.062 \\
ROSMA / Wire Chaser I & 0.483 & 0.060 \\
\end{longtable}
}

The observed diagnostic was regenerated and its aggregate outputs agreed
within 1e-12 with the first run. Code is
\texttt{scripts/run\_matched\_identity\_diagnostic.py}; aggregate
results and the input/output hash manifest are in the dated correction
directory. Individual predictions remain private. Reproduction requires
restoring the authorized private feature and null-prediction inputs and
running this diagnostic before rebuilding the correction and submission
package.

\bibliographystyle{unsrtnat}
\bibliography{references}
\end{document}
